\documentclass[11pt]{article}

% ============================================================
%  Delta m_d (B_d - B_dbar mixing) in warped (Randall--Sundrum) models
%  A slow, pedagogical internal review.
%  Written to track (a) the physics, (b) what is in OUR code,
%  (c) the subtleties and gaps.
%  This is an internal working note, not for distribution.
% ============================================================

\usepackage[margin=1in]{geometry}
\usepackage{amsmath,amssymb,amsthm}
\usepackage{mathtools}
\usepackage{graphicx}
\usepackage{booktabs}
\usepackage{enumitem}
\usepackage{xcolor}
\usepackage{framed}
\usepackage[colorlinks=true,linkcolor=blue,citecolor=blue,urlcolor=blue]{hyperref}

% lightweight "callout" box
\colorlet{shadecolor}{gray!12}
\newenvironment{callout}[1]{%
  \par\medskip\begin{shaded}\noindent\textbf{#1}\par\smallskip}{\end{shaded}\medskip}

\newcommand{\MKK}{M_{\mathrm{KK}}}
\newcommand{\LIR}{\Lambda_{\mathrm{IR}}}
\newcommand{\sw}{s_W}
\newcommand{\cw}{c_W}
\newcommand{\half}{\tfrac{1}{2}}
\newcommand{\vev}{v}
\newcommand{\eps}{\varepsilon}
\newcommand{\Mtwelve}{M_{12}}
\newcommand{\Bd}{B_d}
\newcommand{\Bdbar}{\bar B_d}
\newcommand{\dmd}{\Delta m_d}
\newcommand{\epsK}{\varepsilon_K}
\newcommand{\code}[1]{\texttt{#1}}
\newcommand{\todo}[1]{{\color{red}[#1]}}

\theoremstyle{definition}
\newtheorem{definition}{Definition}

\title{\bf $B_d$--$\bar B_d$ Mixing: $\Delta m_d$ in Warped Models\\[4pt]
\large A slow, pedagogical internal review --- physics, code, and subtleties}
\author{Internal working note (5D--Neutrino--Mixing repo)}
\date{\today}

\begin{document}
\maketitle

\begin{abstract}
\noindent
This note is a from-scratch, deliberately slow walk-through of the neutral
$B_d$-meson mass splitting $\dmd$ and how it constrains a Randall--Sundrum (RS)
warped extra dimension. It assumes no prior knowledge of meson mixing, the
$\Delta F=2$ effective Hamiltonian, or the warped flavour mechanism: all are
built up from the beginning. The central message is a \emph{contrast}: $\dmd$ is
a CP-\emph{conserving} magnitude --- the real part of an off-diagonal mass-matrix
element --- and is therefore markedly \emph{less} constraining than its CP-odd
cousin $\epsK$, which feeds on an imaginary part with a large chiral
enhancement. Part~\ref{part:sm} defines the observable and its Standard-Model
(SM) reference. Part~\ref{part:rs} explains how RS generates a tree-level
Kaluza--Klein (KK) gluon contribution, where the left--right (LR) operators come
from, and where the few-TeV floor sits; it contains a fully worked numerical
example. Part~\ref{part:custodial} asks whether custodial symmetry or flavour
alignment moves the bound (short answer: not really, for this observable).
Part~\ref{part:code} documents exactly what our code does
(\code{quarkConstraints/deltaf2.py}, the \code{physics\_adapters/deltaf2.py}
boundary, and the catalogue constraint \code{B001}, which is
\code{Severity.HARD} and \emph{rigorous}). Part~\ref{part:subtle} catalogues the
conventions and gaps. An annotated reading guide closes the note. The shared
$\Delta F=2$ machinery (operator basis, KK-gluon matching, RG running, matrix
elements) is the \emph{same} core that drives the other meson systems; we state
it self-containedly here, in the form the $B_d$ evaluation actually uses, rather
than relying on an external companion.
\end{abstract}

\tableofcontents

% ====================================================================
\part*{Orientation: why this note exists}
\addcontentsline{toc}{section}{Orientation: why this note exists}
% ====================================================================

Our constraint catalogue contains four neutral-meson mixing systems --- the
kaon ($\epsK$ and $\Delta m_K$), $B_d$, $B_s$, and $D^0$ --- all computed by a
\emph{single} $\Delta F=2$ core (\code{quarkConstraints/deltaf2.py}). They share
the operators, the hadronic-matrix-element structure, and the QCD running. What
distinguishes them is (i) which quark bilinear they probe, (ii) how big the SM
amplitude is, and (iii) --- crucially --- whether the constraint reads the
\emph{magnitude} or the \emph{phase} of the new-physics amplitude. The kaon's
$\epsK$ reads a phase; $\dmd$ reads a magnitude. That one structural difference
is why $\epsK$ pushes the warped KK scale to $\sim$tens of TeV while $\dmd$
typically settles for a few TeV. This note exists to make that difference
precise --- both as physics and as code --- because in our catalogue
\code{B001} ($\dmd$) is a \code{Severity.HARD}, \emph{rigorous} veto, and we want
to know exactly what it is and is not doing.

A one-paragraph executive summary, to be unpacked over the following pages:

\begin{quote}
\small
Two neutral $B_d$ states, $\Bd=(\bar b d)$ and $\Bdbar=(b \bar d)$, are mixed by a
$\Delta B=2$ amplitude $\Mtwelve$; the observed oscillation frequency is the mass
splitting $\dmd = 2|\Mtwelve|$. In RS, light quarks live near the UV brane and the
heavy KK gluons near the IR brane, so a tree-level KK-gluon exchange generates new
$\Delta F=2$ four-quark operators suppressed by $1/\MKK^2$, including dangerous
left--right (LR) operators absent in the SM. Because $\dmd$ constrains only
$|\Mtwelve^{\rm NP}|$ --- a \emph{magnitude} --- the bound sees the operator
coefficient times an $\mathcal{O}(1)$ chiral factor, with \emph{no} access to the
$\mathcal{O}(1)$ CP phase that $\epsK$ exploits and \emph{no} kaon-sized chiral
blow-up. The net effect: anarchic RS satisfies $\dmd$ for
$\MKK\sim\text{few TeV}$, an order of magnitude below the $\sim20$~TeV that $\epsK$
demands. Our code vetoes a point when $|\Mtwelve^{\rm NP}|$ exceeds an
uncertainty-aware ``SM-versus-experiment room'' budget, $\sim\dmd^{\rm exp}/2$ at
leading order.
\end{quote}

\noindent The logical spine of the whole note, in one picture (every symbol ---
$\Mtwelve$, the VLL/VRR/LR operators $C_4/C_5$, the chiral factor $r_\chi$, the
custodial $P_{LR}$, $g_s^*$ --- is defined where it is first used below; read this
as a map to return to, not a self-contained statement):
\begin{footnotesize}
\begin{verbatim}
  B_d = (bbar d),  B_dbar = (b dbar)   --mixed by-->   M_12  (Delta B = 2 amplitude)
                          |
   observable: Delta m_d = 2 |M_12|   ==  a MAGNITUDE (real part), not a phase
                          |
   RS: light quarks UV, KK gluon IR  ==>  tree KK-gluon exchange  ~ (g_s*^2 / M_KK^2)
                          |              generates VLL, VRR, and LR (C4,C5) operators
                          |
   LR operators chirally + RG enhanced  --but for B_d only mildly (r_chi ~ 1.6)
                          |                  [kaon: bare r_chi ~ 26 x running ~8 = ~200,
                          |                   hence eps_K is the killer]
                          |
   CONSTRAINT: |M_12^NP| < budget (uncertainty-aware band)  ==  HARD veto (B001)
                          |
   eps_K (imag part, chiral x running ~200)  ==> ~20 TeV  <==> Delta m_d (real) ==> ~few TeV
                          |
   custodial P_LR: protects Z b_L b_L, NOT the KK-gluon Delta m_d  ==> B001 unchanged
\end{verbatim}
\end{footnotesize}
\noindent Keep this map in mind; every box in it is unpacked below. The single
most common confusion --- that $\dmd$ and $\epsK$ are ``the same constraint in two
mesons'' --- is already half-defused here: they are the \emph{real} and
\emph{imaginary} projections of the same kind of amplitude, and those projections
carry wildly different enhancement factors.

% ====================================================================
\part{The physics from scratch: what $\Delta m_d$ is}
\label{part:sm}
% ====================================================================

\section{Mixing: two states, one amplitude}
\label{sec:mixing}

A neutral $B_d$ meson is a bound state $\Bd=(\bar b\, d)$ with bottom number
$+1$; its antiparticle is $\Bdbar=(b\,\bar d)$. The strong and electromagnetic
interactions conserve flavour, so naively these two are distinct, stable
flavour eigenstates. The weak interaction does not conserve $b$- and $d$-number
separately: a second-order weak process can turn a $\Bd$ into a $\Bdbar$. This is
a $\Delta B = 2$ transition (bottom number changes by two units), and it makes
the two flavour states mix.

The two-state system is governed by an effective $2\times2$ Hamiltonian
$H = M - \tfrac{i}{2}\Gamma$, where $M$ (mass matrix) and $\Gamma$ (decay matrix)
are Hermitian. The off-diagonal entry $\Mtwelve \equiv M_{12}$ is the dispersive
(virtual) mixing amplitude; $\Gamma_{12}$ is the absorptive (on-shell) one. The
physical (mass) eigenstates $B_{d,H}$ and $B_{d,L}$ (heavy/light) are linear combinations of
$\Bd$ and $\Bdbar$, and they are split in mass by
\begin{equation}
\boxed{\;\dmd \;\equiv\; m_{B_H} - m_{B_L} \;\simeq\; 2\,|\Mtwelve|\;,\;}
\label{eq:dmddef}
\end{equation}
the approximation holding to excellent accuracy because $|\Gamma_{12}|\ll
|\Mtwelve|$ in the $B_d$ system. Experimentally, $\dmd$ is measured as an
\emph{oscillation frequency}: a meson born as $\Bd$ oscillates into $\Bdbar$ and
back with angular frequency $\dmd/\hbar$. The companion observables stored
alongside it are the dimensionless ratio $x_d = \dmd/\Gamma_d$ and the
time-integrated mixing probability $\chi_d$.

\begin{definition}[The observable]
$\dmd$ is the mass difference of the two neutral-$B_d$ mass eigenstates, equal to
twice the magnitude of the dispersive mixing amplitude $\Mtwelve$. It is
\textbf{CP-conserving}: it depends on $|\Mtwelve|$, not on $\arg\Mtwelve$. The
CP-violating information lives in the \emph{phase} of $\Mtwelve$ (relative to the
decay amplitudes), which is a separate observable.
\end{definition}

This is the conceptual fork in the road, and it is worth pausing on. Write
$\Mtwelve = |\Mtwelve|\,e^{i\phi}$. Then $\dmd$ sees $|\Mtwelve|$ and is blind to
$\phi$; the CP asymmetries ($\sin 2\beta$, $\phi_s$ in the $B_s$ system, $\epsK$
in the kaon system) see $\phi$ and are blind to a real shift of $|\Mtwelve|$. A
new-physics contribution that is \emph{aligned} in phase with the SM adds to
$|\Mtwelve|$ but generates no new CP violation; one that is \emph{misaligned}
moves the phase. The kaon's $\epsK$ is so powerful precisely because it isolates
an imaginary part that is tiny in the SM, so any new imaginary piece stands out.
$\dmd$, by contrast, sits on top of a large, real SM amplitude and can only
constrain a \emph{shift} of a number that is already $\mathcal{O}(1)$ in SM
units. We return to this contrast quantitatively in
\S\ref{sec:realvsimag}.

\section{The SM expectation and the measurement}
\label{sec:smvalue}

In the SM, $\Mtwelve$ comes from the famous box diagrams: two $W$ bosons and two
up-type quarks (dominantly the top) running in a loop, with external $b$ and $d$
legs. The amplitude is therefore CKM- and loop-suppressed, scaling as
$\Mtwelve^{\rm SM}\propto (V_{tb}V_{td}^*)^2\, S_0(x_t)\, f_{B_d}^2\hat B_{B_d}$,
where $S_0$ is the Inami--Lim loop function, $f_{B_d}$ the decay constant, and
$\hat B_{B_d}$ the bag parameter. The CKM combination $V_{tb}V_{td}^*\sim
V_{td}$ is what makes $\dmd$ a probe of the $1\!\leftrightarrow\!3$ down-sector
transition.

The numbers our catalogue uses (sidecar \code{B001.yaml}) are, for the
\emph{experiment}, the HFLAV/PDG~2025 world average
\begin{equation}
\dmd^{\rm exp} = 0.5069 \pm 0.0019~\text{ps}^{-1},
\label{eq:dmdexp}
\end{equation}
and, for the \emph{theory}, the HPQCD~2019 lattice prediction (arXiv:1907.01025)
\begin{equation}
\dmd^{\rm SM} = 0.555\,{}^{+0.028}_{-0.055}\,(0.029)~\text{ps}^{-1}
   \;=\; 0.555 \pm 0.062~\text{ps}^{-1},
\label{eq:dmdsm}
\end{equation}
where the symmetrised uncertainty combines the larger asymmetric side
($0.055$) with the second uncertainty ($0.029$) in quadrature,
$\sqrt{0.055^2+0.029^2}=0.062$, exactly the \code{uncertainty\_policy} recorded
in the YAML. The two agree at about $0.8\sigma$ --- the SM already accounts for
the entire measured oscillation. Note the catalogue stores both the modern
Belle~II value ($0.516\pm0.008\pm0.005$~ps$^{-1}$, arXiv:2302.12791) under
\code{recent\_experimental\_inputs} and the legacy Csaki--Falkowski--Weiler RS
context bounds ($21$ and $33$~TeV, arXiv:0804.1954) under
\code{paper\_era\_reference}; only the HFLAV/PDG average and the HPQCD SM number
feed the live constraint.

The fact that SM and experiment agree is what \emph{defines the room} for new
physics: any RS contribution to $|\Mtwelve|$ must fit inside the gap between them,
dressed by uncertainties. That ``room'' is the budget of \S\ref{sec:budget}.

\section{The $\Delta F=2$ effective operators}
\label{sec:operators}

Below the KK scale, the new physics is invisible as particles and shows up only
as local four-quark operators --- the language of effective field theory. For a
$\bar b d \to b \bar d$ ($\Delta B=2$) transition, the complete dimension-six
basis has eight operators; the subset our tree-level KK-gluon matching generates
is
\begin{align}
Q_1^{\rm VLL} &= (\bar d_L\gamma_\mu b_L)(\bar d_L\gamma^\mu b_L), &
Q_1^{\rm VRR} &= (\bar d_R\gamma_\mu b_R)(\bar d_R\gamma^\mu b_R),
\label{eq:VLL}\\
Q_4^{\rm LR} &= (\bar d_R\, b_L)(\bar d_L\, b_R), &
Q_5^{\rm LR} &= (\bar d_R^\alpha b_L^\beta)(\bar d_L^\beta b_R^\alpha),
\label{eq:LR}
\end{align}
with $\alpha,\beta$ colour indices. The VLL operator \eqref{eq:VLL} has the same
chiral structure as the SM box; VRR is its parity mirror; the two LR operators
\eqref{eq:LR} are \emph{chirality-mixing} (scalar/tensor-like) and have no SM
counterpart. (The full basis also contains SLL/SRR scalar operators, which a pure
vector KK-gluon exchange does not generate at tree level; see
\S\ref{sec:gap}.) The effective Hamiltonian is
\begin{equation}
\mathcal{H}^{\Delta B=2}_{\rm eff}
   = \sum_i C_i(\mu)\, Q_i(\mu),
\label{eq:Heff}
\end{equation}
and the mixing amplitude is the meson matrix element
\begin{equation}
\Mtwelve^{\rm NP}
   = \langle \Bdbar |\, \mathcal{H}^{\Delta B=2}_{\rm eff}\, | \Bd\rangle
   = \sum_i C_i(\mu)\, \langle \Bdbar | Q_i(\mu) | \Bd\rangle .
\label{eq:M12sum}
\end{equation}
Two features will dominate everything downstream:
\begin{itemize}[leftmargin=1.4em]
\item \textbf{The LR operators are special.} Their hadronic matrix elements are
enhanced relative to the VLL one by a \emph{chiral factor}
$r_\chi = \big(m_P/(m_{q_1}+m_{q_2})\big)^2$ (for $B_d$,
$m_P=m_{B_d}$, $m_{q_1}+m_{q_2}=m_b+m_d$), which is large because the meson mass
is much bigger than the sum of its constituent current-quark masses --- the
pseudoscalar density is fixed by the equations of motion to $\sim m_P^2/(m_{q_1}+m_{q_2})$.
They are also enhanced
by QCD running between the matching scale and the hadronic scale (the LR
anomalous dimension is large). This is the same mechanism that makes LR operators
the dominant danger in \emph{every} meson system.
\item \textbf{But the $B_d$ chiral factor is modest.} For $B_d$,
$r_\chi=(5.28/(4.18+0.005))^2\approx 1.6$ --- order unity. For the kaon, with the
repo's FLAG $2$~GeV light-quark masses the \emph{bare} chiral factor is
$r_\chi=(m_K/(m_s+m_d))^2\approx 26$, and this is multiplied by the LR QCD-running
factor $\sim8$ (the $C_4$--$C_5$ pair mixes under a $2\times2$ anomalous-dimension
matrix that enhances the LR contribution between the matching and hadronic
scales). The two combine to a kaon LR enhancement of order
$26\times8\sim200$ relative to the VLL one
(the breakdown is worked through in the kaon note's running section; note that
Csaki--Falkowski--Weiler quote a \emph{combined} chiral-plus-running factor
$\tfrac34\,(m_K/(m_s+m_d))^2\,\eta_1^{-5}\approx18$ in their normalisation, which
already folds in the running and so should not be multiplied by it again). For
$B_d$ the running enhancement is similar in kind but the chiral prefactor is
$\sim1.6$, not $\sim26$, so the $B_d$ LR operator ends up only mildly enhanced.
\textbf{The $B_d$ LR operator is enhanced, but nowhere near as violently as the
kaon's.} This is half the reason $\dmd$ is weaker than $\epsK$; the other half is
the magnitude versus phase distinction.
\end{itemize}

% ====================================================================
\part{How RS generates $\Delta m_d$}
\label{part:rs}
% ====================================================================

\section{The warped flavour mechanism, in just enough detail}
\label{sec:rsmech}

Recall the RS picture (consistent with the repo conventions). Spacetime is a
slice of AdS$_5$ between a UV brane (Planck scale) and an IR brane (TeV scale).
Bulk fermions carry a dimensionless bulk-mass parameter $c$ that localises their
zero mode: $c>\tfrac12$ peaks the wavefunction toward the UV brane (light
fermions), $c<\tfrac12$ toward the IR brane (heavy fermions). The Higgs sits on
the IR brane, so a fermion's mass is set by how much of its wavefunction reaches
the IR --- this is the geometric origin of the Yukawa hierarchy, with
$\mathcal{O}(1)$ (\emph{anarchic}) 5D couplings.

Now bring in the gauge sector. The gluon has a KK tower; the first KK gluon is
localised near the IR brane and is heavy, $\MKK\sim$~few TeV. Its coupling to a
given quark depends on the overlap of the (flat-ish) zero-mode quark profile and
the (IR-peaked) KK-gluon profile. Two consequences follow, and both matter for
$\dmd$:
\begin{enumerate}[leftmargin=1.6em]
\item \textbf{The KK gluon couples non-universally.} IR-localised (heavy) quarks
couple to the KK gluon much more strongly than UV-localised (light) ones. The
coupling is enhanced over the ordinary 4D strong coupling by a volume factor
$\sqrt{2k\pi r_c}\sim\sqrt{2L}$; this is the $g_s^*\sim 6$ at tree level
documented in \code{quarkConstraints/couplings.py} (citing
Csaki--Falkowski--Weiler 0804.1954).
\item \textbf{Non-universal $\Rightarrow$ flavour-violating.} Because the coupling
is generation-dependent, rotating to the quark mass basis does \emph{not}
diagonalise the KK-gluon coupling. Off-diagonal couplings survive --- a
tree-level flavour-changing neutral current. This is the warped analogue of an
``RS-GIM'' partial protection: the off-diagonal couplings are suppressed by the
small UV-localised light-quark wavefunctions, but not zero.
\end{enumerate}
Integrating out the first KK gluon at tree level produces exactly the
$\Delta F=2$ operators of \eqref{eq:VLL}--\eqref{eq:LR}, with coefficients set by
the off-diagonal mass-basis couplings and a $1/\MKK^2$ propagator.

\section{Parametric size and where the floor comes from}
\label{sec:parametric}

Schematically, the Wilson coefficient of a $\Delta B=2$ operator from KK-gluon
exchange is
\begin{equation}
C_i \;\sim\; \frac{(g_s^*)^2}{\MKK^2}\; (\Delta_L)_{db}\,(\Delta_R)_{db}\,,
\label{eq:Cscale}
\end{equation}
where $(\Delta_{L,R})_{db}$ are the left/right off-diagonal $d$--$b$ couplings in
the mass basis. For the chirality-mixing LR operators, \emph{both} a left and a
right coupling appear; for VLL only left couplings, for VRR only right. The
mixing amplitude is then $\Mtwelve^{\rm NP}\sim C_i\langle Q_i\rangle$, and the
constraint
\begin{equation}
|\Mtwelve^{\rm NP}| \;\lesssim\; \tfrac12\,\dmd^{\rm exp}
\label{eq:floorcondition}
\end{equation}
turns into a \emph{lower bound on $\MKK$}, because $|\Mtwelve^{\rm NP}|\propto
1/\MKK^2$. Plugging in anarchic $\mathcal{O}(1)$ couplings dressed by the warped
suppression of the light-quark profiles, the $B_d$ bound from the LR operator
comes out at the few-TeV level for the chirality-flip coefficient. This is the
floor.

It is illuminating to put the four meson systems side by side, because they all
use \emph{the same} formula \eqref{eq:Cscale} and differ only in the flavour
indices, the matrix elements, and --- decisively --- whether the constraint reads
a magnitude or a phase:
\begin{center}
\begin{tabular}{llll}
\toprule
System & Transition & Constraint reads & Anarchic $\MKK$ floor \\
\midrule
$\epsK$ & $s\!\leftrightarrow\! d$ & \textbf{Im}$\Mtwelve$ (phase) & $\sim 20$~TeV \\
$\dmd$ (\textbf{this note}) & $b\!\leftrightarrow\! d$, $\sim V_{td}$ & $|\Mtwelve|$ (magnitude) & $\sim$ few TeV \\
$\Delta m_s$ / $\phi_s$ & $b\!\leftrightarrow\! s$, $\sim V_{ts}$ & magnitude / phase & $\sim 2$--$3$~TeV \\
$D^0$ mixing & $c\!\leftrightarrow\! u$ (up) & magnitude & $\sim$ few TeV \\
\bottomrule
\end{tabular}
\end{center}
\noindent $D^0$ is also only a few-TeV constraint even though, like $\dmd$, it
reads a magnitude: it is an up-sector ($c$--$u$) transition, so it carries neither
the large down-sector chiral enhancement nor a kaon-sized CKM/SM hierarchy to
contrast against. Magnitude-reading down-sector ($\dmd$, $\Delta m_s$) and the
up-sector ($D^0$) therefore land in the same few-TeV band, while only the
phase-reading $\epsK$ climbs to tens of TeV.
The $\sim20$~TeV figure for $\epsK$ is the value the Csaki--Falkowski--Weiler and
Blanke et al.\ analyses (0804.1954, 0809.1073) quote for the lightest KK gluon in
the fully anarchic scenario; $\dmd$ lives an order of magnitude below it. That
gap --- same model, same machinery, factor of $\sim$$7$--$10$ in the scale --- is
the entire reason this note exists.

\section{Real versus imaginary: the quantitative contrast}
\label{sec:realvsimag}

Why is reading a magnitude so much weaker than reading a phase? Three multiplying
factors, all in $\epsK$'s favour:
\begin{enumerate}[leftmargin=1.8em]
\item \textbf{The SM imaginary part is tiny.} $\epsK$ measures $\mathrm{Im}
\Mtwelve^{K}$, which in the SM is suppressed by a small CKM phase. The NP
contribution competes against a small number, so a modest NP imaginary piece is a
large fractional perturbation. $\dmd$ measures $|\Mtwelve^{B_d}|$, which sits on
top of the \emph{full} SM amplitude --- a large, well-measured number. NP can only
nudge it.
\item \textbf{The anarchic phase is generic $\mathcal{O}(1)$.} In anarchic RS the
NP contribution carries an uncorrelated $\mathcal{O}(1)$ CP phase, so the
imaginary part is not suppressed on the NP side --- $\epsK$ gets the full
$|\Mtwelve^{\rm NP}|$ times $\sin(\mathcal{O}(1))$. The phase is ``free'' danger;
the magnitude has no such bonus.
\item \textbf{The kaon chiral enhancement is huge; the $B_d$ one is mild.} As
above, the kaon LR operator is enhanced by $\sim$$200$ = (bare chiral $r_\chi\approx26$)
$\times$ (LR running $\sim8$) relative to the VLL operator, whereas the $B_d$
chiral prefactor is only $r_\chi\approx1.6$. So even operator-for-operator, the
kaon amplitude is amplified far more.
\end{enumerate}
Multiply these and one understands the $\sim20$~TeV (kaon) versus few-TeV ($B_d$)
hierarchy without a single Feynman diagram. \textbf{Slogan:} \emph{$\epsK$ wins on
three counts --- tiny SM imaginary part, free $\mathcal{O}(1)$ NP phase, and a
giant chiral enhancement --- and $\dmd$ has none of them; it is the same physics
read through a much blunter projection.}

\begin{callout}{Worked example --- the $B_d$ budget and a KK-scale estimate.}
\textbf{Step 1: the matrix elements.} With the in-code constants
$f_{B_d}=0.1900$~GeV, $m_{B_d}=5.27972$~GeV, $m_b=4.18$~GeV,
$m_d=0.00467$~GeV, and bag parameters $B_1=0.87,\ B_4=1.02,\ B_5=0.96$, the chiral
ratio is $r_\chi=(m_{B_d}/(m_b+m_d))^2\approx1.59$ and
\[
\langle Q_1^{\rm VLL}\rangle = \tfrac23 f_{B_d}^2 m_{B_d} B_1 \approx 0.111~\text{GeV}^3,
\quad
\langle Q_4^{\rm LR}\rangle \approx 0.100~\text{GeV}^3,
\quad
\langle Q_5^{\rm LR}\rangle \approx 0.161~\text{GeV}^3.
\]
Note the LR matrix element is comparable to (not the $\sim$$26\times$ of the
kaon's chirally enhanced matrix element) the VLL one --- the
$B_d$ chiral enhancement is mild, exactly as advertised.

\textbf{Step 2: the budget.} The experimental world average
\eqref{eq:dmdexp} converts to GeV via the in-code anchor
$\dmd^{\rm exp}=3.334\times10^{-13}$~GeV (so
$1~\text{ps}^{-1}\leftrightarrow6.58\times10^{-13}$~GeV). The SM central from the
core convention is $\dmd^{\rm SM}=3.6\times10^{-13}$~GeV. The
\emph{leading-order} core budget is
\[
\text{budget}_{\rm core}=\max\!\Big(\tfrac{\dmd^{\rm exp}}{2},\
\tfrac{|\dmd^{\rm exp}-\dmd^{\rm SM}|}{2}\Big)
=\tfrac{3.334\times10^{-13}}{2}=1.667\times10^{-13}~\text{GeV}.
\]
The catalogue \emph{HARD} budget instead uses the uncertainty-aware room
\[
\text{budget}_{\rm B001}=\frac{|\dmd^{\rm exp}-\dmd^{\rm SM}|+\sqrt{\sigma_{\rm exp}^2+\sigma_{\rm SM}^2}}{2}
\approx 3.38\times10^{-14}~\text{GeV},
\]
with $\sigma_{\rm SM}$ dominating ($\sigma_{\rm SM}\approx4.1\times10^{-14}$,
$\sigma_{\rm exp}\approx1.2\times10^{-15}$~GeV). The central (no-$\sigma$) room is
$1.33\times10^{-14}$~GeV. \textbf{So the live veto is the loose band,
$\sim3.4\times10^{-14}$~GeV} --- five times tighter than the naive
$\dmd^{\rm exp}/2$, because it folds in the SM/exp agreement.

\textbf{Step 3: invert for the KK scale (the headline ``few TeV'').} Now close the
loop. The dominant operator is $Q_4^{\rm LR}$, with coefficient
$C_4^{\rm LR}=-(g_s^*)^2(\Delta_L)_{db}(\Delta_R)_{db}/\MKK^2$
[Eq.~\eqref{eq:Cscale} with the $-1$ colour factor of Eq.~\eqref{eq:codeLR}], so
\[
|\Mtwelve^{\rm NP}|\;\approx\;\frac{(g_s^*)^2\,|(\Delta_L)_{db}(\Delta_R)_{db}|}{\MKK^2}\,
\langle Q_4^{\rm LR}\rangle .
\]
For the inputs take the volume-enhanced strong coupling $g_s^*\approx6$
(\S\ref{sec:rsmech}, so $(g_s^*)^2\approx36$) and the \emph{anarchic} RS-GIM
off-diagonal couplings. These are not free: RS-GIM ties them to the same overlap
factors that set the light-quark masses, $(\Delta_{L,R})_{db}\sim f_{d}\,f_{b}$,
which for the $d$--$b$ entry land near $(\Delta_L)_{db}\sim(\Delta_R)_{db}\sim
3\times10^{-4}$ in units where the diagonal heavy coupling is $\mathcal{O}(1)$
(the geometric suppression of the UV-localised light-quark profiles). Then
$|(\Delta_L)_{db}(\Delta_R)_{db}|\sim9\times10^{-8}$ and
\[
|\Mtwelve^{\rm NP}|\;\approx\;\frac{36\times(9\times10^{-8})}{(\MKK/\text{GeV})^2}\times0.10~\text{GeV}^3
\;=\;\frac{3.2\times10^{-7}~\text{GeV}^3}{(\MKK/\text{GeV})^2}\,\text{GeV}.
\]
Setting this equal to the live budget $\approx3.4\times10^{-14}$~GeV and solving,
\[
\MKK\;\approx\;\sqrt{\frac{3.2\times10^{-7}}{3.4\times10^{-14}}}~\text{GeV}
\;\approx\;3\times10^{3}~\text{GeV}\;\approx\;\textbf{3 TeV},
\]
which is the order-of-magnitude origin of the few-TeV floor that the whole note
rests on. The result is only logarithmically sensitive to the inputs (it depends
on $\sqrt{(g_s^*)^2|\Delta\Delta|\langle Q_4\rangle/\text{budget}}$), so the
residual $\mathcal{O}(1)$ anarchic spread and the running of $C_4$ move it by
factors of order one, not orders of magnitude. (Contrast the kaon: the same
inversion run on $\mathrm{Im}\,C_4^K$ --- which competes against a tiny SM
imaginary part rather than the full amplitude, and carries the $\sim200\times$
chiral-plus-running enhanced kaon LR contribution --- pushes the floor an order of magnitude higher,
to $\sim20$~TeV. The three multiplying advantages of $\epsK$ over $\dmd$ are
itemised in \S\ref{sec:realvsimag}.)

\textbf{Step 4: contrast with $\epsK$.} The kaon budget is
$|\epsK^{\rm exp}-\epsK^{\rm SM}|=|2.228-2.161|\times10^{-3}=6.7\times10^{-5}$
(dimensionless), read off the \emph{imaginary} part of $\Mtwelve^K$ through a
$1/\Delta m_K$ amplification. The kaon constraint is sharper not because its
budget is smaller in some common unit, but because (i) it taps the small SM
imaginary part, (ii) it gets the free anarchic phase, and (iii) its LR
contribution (chiral $\times$ running) is $\sim$$200\times$ bigger. The single benchmark above contains the whole
story: for the same KK scale, $|M_{12}^{B_d,\,\rm NP}|$ must merely stay below a
number of order the full $B_d$ amplitude, whereas $\mathrm{Im}\,M_{12}^{K,\,\rm NP}$
must stay below a number two-plus orders of magnitude below the full kaon
amplitude.
\end{callout}

% ====================================================================
\part{Model, custodial, and alignment dependence}
\label{part:custodial}
% ====================================================================

\section{Does custodial symmetry change this bound?}
\label{sec:custodial}

The custodial-symmetry story that dominates the electroweak sector (the
$SU(2)_L$-type bulk gauge enlargement to $SU(2)_L\times SU(2)_R\times U(1)_X$ with a
discrete $P_{LR}$) is, for $\dmd$, almost a non-event --- and this is worth saying
clearly because it is a common source of confusion. Custodial protection does two
specific things: (a) the continuous $SU(2)_R$ kills the volume-enhanced $T$
parameter, and (b) the discrete $P_{LR}$ protects the $Z\bar b_L b_L$ vertex (and,
in the flavour sector, the \emph{tree-level $Z$-mediated} flavour-changing
coupling to left-handed down quarks).

But $\dmd$ in our code is generated by \textbf{KK-gluon} exchange, not by
$Z$ exchange. The $P_{LR}$ protection of the $Zd_Ld_L$ coupling does \emph{not}
touch the gluonic operators, because the gluon is a colour octet and is blind to
the $SU(2)_R$/$P_{LR}$ structure that protects the $Z$ couplings. Blanke et al.\
(0809.1073) make exactly this point: $P_{LR}$ strongly suppresses the
$Z$-mediated $\Delta F=2$ contribution to left-handed down quarks, but this
\emph{does not directly relax the $B_d$ mixing bound} from the KK gluon --- it
rather removes a competing $Z$ contribution and lets the gluonic one stand. So
the headline is:
\begin{center}
\fbox{\parbox{0.92\linewidth}{\centering
\textbf{Custodial $P_{LR}$ does \emph{not} move the $\dmd$ / \code{B001} bound.}
It protects $Z\bar b_L b_L$ (handled separately in our \code{T010}/\code{T011}),
not the KK-gluon $\Delta B=2$ operators that drive $|M_{12}^{B_d,\,\rm NP}|$.}}
\end{center}
This is the concrete meaning of the brief's tag ``B001 is not affected by
custodial.'' In our paired minimal-vs-custodial scans, turning on the custodial
mode removes the $Z\to b\bar b$ veto but leaves the $\Delta F=2$ set
($\epsK$/$\Delta m_s$/$\phi_s$, with $\dmd$ alongside) as the rigorous floor.

\section{Does flavour alignment change it?}
\label{sec:alignment}

Yes --- and this is the lever that actually matters for $\dmd$. The bound
\eqref{eq:floorcondition} assumes \emph{anarchic} flavour: $\mathcal{O}(1)$,
uncorrelated 5D Yukawas with no special structure. If instead one
\emph{aligns} the bulk masses / 5D Yukawas so that the off-diagonal mass-basis
couplings $(\Delta_{L,R})_{db}$ are suppressed beyond their generic anarchic
size, the Wilson coefficient \eqref{eq:Cscale} shrinks and the KK floor drops.
Blanke et al.\ (0809.1073) exhibit exactly such regions: ``only modest
fine-tuning'' suffices to satisfy all $\Delta F=2$ and electroweak constraints
down to $\MKK\simeq3$~TeV. Alignment is a statement about the \emph{flavour}
structure of the 5D couplings, orthogonal to the custodial \emph{gauge}
structure; the two evade different bounds (alignment for $\Delta F=2$, custodial
for electroweak).

\textbf{What our code does about this:} it does \emph{not} impose anarchy or
alignment by fiat --- it computes $(\Delta_{L,R})_{db}$ from the actual fitted
mass-basis couplings of the scanned point (the \code{quark\_mass\_basis\_couplings}
extra; \S\ref{sec:codeflow}). A point that happens to be aligned will simply
produce a small $|\Mtwelve^{\rm NP}|$ and pass; a generic anarchic point will
produce a large one and fail. The constraint is therefore \emph{honest about
alignment}: it rewards genuinely aligned points and penalises generic ones,
rather than baking in an assumption.

\section{Rigorous or proxy?}
\label{sec:rigor}

For $\dmd$, our treatment is \textbf{rigorous}, not a proxy, and it is worth
being precise about what that means here (contrast the oblique $S,T,U$ sector,
which is an avowed proxy). The $\dmd$ path:
\begin{itemize}[leftmargin=1.4em]
\item uses the \emph{actual} fitted mass-basis couplings of the point, not a
parametric estimate;
\item builds genuine Wilson coefficients from a tree-level KK-gluon matching;
\item QCD-runs them from the matching scale to $\mu_{\rm had}=2$~GeV with a real
RG evolution;
\item contracts them with FLAG-2024 hadronic matrix elements (the
$f_{B_d},B_i$ above); and
\item compares $|\Mtwelve^{\rm NP}|$ to an uncertainty-aware budget built from the
HFLAV/PDG experiment and the HPQCD SM number.
\end{itemize}
The honest caveats (scalar SLL/SRR operators not generated, single-KK-mode
truncation, leading-log running, GeV/ps anchor) are conventions and deferred
refinements, \emph{not} parametric guesses; they are catalogued in
\S\ref{part:subtle}. By the catalogue's own taxonomy this is a \code{rigorous}
constraint with documented approximations, which is why \code{B001} carries
\code{Severity.HARD}.

% ====================================================================
\part{What is in OUR code}
\label{part:code}
% ====================================================================

\emph{Reader here for the physics may skim Part~\ref{part:code} on a first pass
and jump to the subtleties in Part~\ref{part:subtle}; this part documents the
exact code so the note doubles as an implementation reference. The one physical
takeaway to retain from all of it: the code literally calls \code{abs()} on the
complex $\Mtwelve^{\rm NP}$ before comparing to the budget (\S\ref{sec:codeME}) ---
that single line is what makes $\dmd$ a magnitude constraint rather than a phase
one, and it is the code-level shadow of the whole real-vs-imaginary thesis.}

\section{The three layers}
\label{sec:layers}

The $\dmd$ evaluation passes through three files, in order:
\begin{enumerate}[leftmargin=1.6em]
\item \code{quarkConstraints/deltaf2.py} --- the shared $\Delta F=2$ \emph{core}:
operator matching, hadronic matrix elements, QCD running, and the per-system
evaluators.
\item \code{flavor\_catalog\_constraints/physics\_adapters/deltaf2.py} --- a thin,
append-only \emph{adapter} that exposes a stable surface
(\code{bd\_mixing\_wilsons\_from\_couplings},
\code{bd\_mixing\_from\_wilsons\_with\_running},
\code{bd\_mixing\_core\_inputs}) so the catalogue never imports core internals
directly.
\item \code{flavor\_catalog\_constraints/primary/beauty/B001.py} --- the
\emph{catalogue constraint}: it loads the YAML anchors, builds the
uncertainty-aware budget, calls the adapter, and returns a pass/fail
\code{ConstraintResult}.
\end{enumerate}

\section{The Wilson coefficients, exactly}
\label{sec:codewilsons}

In \code{compute\_delta\_f2\_wilsons} (\code{deltaf2.py}), with
$\text{prefactor}=1/\MKK^2$ and the off-diagonal $d$--$b$ mass-basis couplings
\code{left}, \code{right}, the four coefficients are
\begin{align}
C_1^{\rm VLL} &= \frac{\text{left}\cdot\text{left}}{6\,\MKK^2}, &
C_1^{\rm VRR} &= \frac{\text{right}\cdot\text{right}}{6\,\MKK^2},
\label{eq:codeVLL}\\
C_4^{\rm LR} &= -\,\frac{\text{left}\cdot\text{right}}{\MKK^2}, &
C_5^{\rm LR} &= +\,\frac{\text{left}\cdot\text{right}}{3\,\MKK^2}.
\label{eq:codeLR}
\end{align}
These are matched at $\mu=\MKK$ (the \code{matching\_scale} field). The relative
weights ($1/6$ for the vector colour-singlet projection, $-1$ and $+1/3$ for the
two LR colour structures) are the standard tree-level KK-gluon Fierz/colour
factors. The $B_d$ entry is selected from the default bundle by key \code{"b\_d"}
(generations \code{(0, 2)}, i.e.\ $d$--$b$, sector \code{"down"}), with the
adapter helper \code{\_bd\_wilsons} pulling exactly that entry out of
\code{DEFAULT\_DELTA\_F2\_INPUTS\_V1}.

\section{The matrix elements and $M_{12}$, exactly}
\label{sec:codeME}

The generic per-meson matrix elements live in \code{\_meson\_matrix\_elements}:
with $f_P, m_P$, current-quark masses $m_{q_1},m_{q_2}$, and bag parameters
$B_1,B_4,B_5$,
\begin{align}
\langle Q_1^{\rm VLL}\rangle &= \tfrac23 f_P^2 m_P\, B_1, \\
\langle Q_4^{\rm LR}\rangle &= \big(\tfrac{r_\chi}{6}+\tfrac14\big) f_P^2 m_P\, B_4, &
\langle Q_5^{\rm LR}\rangle &= \big(\tfrac{r_\chi}{2}+\tfrac1{12}\big) f_P^2 m_P\, B_5,
\end{align}
with $r_\chi=(m_P/(m_{q_1}+m_{q_2}))^2$ --- the chiral factor of
\S\ref{sec:operators} made explicit in code. The $B_d$ inputs are the
module-level constants
\begin{equation}
\begin{aligned}
&F\_BD=0.1900~\text{GeV},\ M\_BD=5.27972~\text{GeV},\
M\_B\_QUARK=4.18~\text{GeV},\ M\_D\_QUARK\_BD=0.00467~\text{GeV},\\
&B\_1\_BD=0.87,\quad B\_4\_BD=1.02,\quad B\_5\_BD=0.96
\quad(\text{FLAG 2024}).
\end{aligned}
\end{equation}
The function \code{compute\_m12\_np} forms the coherent sum
\eqref{eq:M12sum},
\begin{equation}
\Mtwelve^{\rm NP}
   = C_1^{\rm VLL}\langle Q_1\rangle + C_1^{\rm VRR}\langle Q_1\rangle
   + C_4^{\rm LR}\langle Q_4\rangle + C_5^{\rm LR}\langle Q_5\rangle,
\label{eq:codeM12}
\end{equation}
(VRR reuses the VLL matrix element by parity) and \code{evaluate\_bd\_mixing}
takes the magnitude \code{abs(m12)}. \textbf{This is the code-level embodiment of
``$\dmd$ is a magnitude'':} the very last step throws away the phase via
\code{abs(\,)}. The phase-preserving companion
\code{bd\_mixing\_m12\_np\_from\_wilsons\_with\_running} exists in the adapter
(returning the complex $\Mtwelve^{\rm NP}$) but is \emph{not} what \code{B001}
uses --- it is there for diagnostics and for the CP observables that \emph{do}
read the phase.

\section{QCD running}
\label{sec:coderunning}

By default the core runs the Wilson coefficients from $\MKK$ down to
$\mu_{\rm had}=2.0$~GeV via leading-log QCD evolution
(\code{\_evolve\_wilsons} $\to$ \code{evolve\_deltaf2\_wilsons}) before
contracting with the $2$-GeV matrix elements. The VLL/VRR coefficients run with
the current-current anomalous dimension; the LR coefficients $C_4,C_5$ run in the
scalar basis matched to the $B_4/B_5$ inputs. \code{B001} always runs
(\code{mu\_had = \_MU\_HAD\_GEV = 2.0}, \code{qcd\_running\_applied: True} in the
diagnostics).

\section{The budget and the veto, exactly}
\label{sec:budget}

One point to settle up front, because it is assembled piecemeal elsewhere in the
note: there are \emph{two} budgets, and they are not the same number. The
``$\text{budget}=\dmd^{\rm exp}/2$'' convention named in the brief and the CLAUDE
notes is the \emph{core/diagnostic} budget \eqref{eq:corebudget}; it is computed
and reported but \textbf{is not what \code{B001} vetoes on}. The live HARD veto
uses the tighter uncertainty-aware band \eqref{eq:b001budget}, about five times
smaller. Whenever this note says ``the budget,'' it means the live band unless
explicitly flagged as the core convention.

This is where \code{B001} differs from the bare core. The core's legacy $B_d$
budget (\code{\_bd\_budget}, mirrored in \code{bd\_mixing\_core\_inputs}) is
\begin{equation}
\text{budget}_{\rm core}
   = \max\!\Big(\tfrac{\dmd^{\rm exp}}{2},\ \tfrac{|\dmd^{\rm exp}-\dmd^{\rm SM}|}{2}\Big)
   = 1.667\times10^{-13}~\text{GeV},
\label{eq:corebudget}
\end{equation}
using the in-code constants \code{DELTA\_M\_BD\_EXP = 3.334e-13} and
\code{DELTA\_M\_BD\_SM = 3.6e-13} (GeV). This is the ``\textbf{budget = experimental
$\Delta m$ over $2$}'' convention referenced in the brief: at leading order the
$B_d$ budget is just $\dmd^{\rm exp}/2$, because the exp/SM gap is smaller than the
experimental value itself. The constraint keeps this as a \emph{diagnostic}
(\code{core\_default\_budget}, \code{core\_default\_ratio}).

The \emph{live} HARD budget that \code{B001} actually vetoes on is the
uncertainty-aware band built in \code{\_build\_budget\_band}:
\begin{equation}
\boxed{\;\text{budget}_{\rm B001}
   = \frac{|\dmd^{\rm exp}-\dmd^{\rm SM}| + \sqrt{\sigma_{\rm exp}^2+\sigma_{\rm SM}^2}}{2}
   \;\approx\; 3.38\times10^{-14}~\text{GeV}\;}
\label{eq:b001budget}
\end{equation}
(the \code{loose\_budget}, assigned to \code{hard\_veto\_budget}). Here the
experiment and its uncertainty come from the YAML
(\code{canonical\_experimental\_average}: $0.5069\pm0.0019$~ps$^{-1}$), converted
to GeV through the anchor \code{DELTA\_M\_BD\_EXP\_GeV} (which the code
\emph{cross-checks} against the core \code{DELTA\_M\_BD\_EXP} to
$10^{-12}$ relative tolerance, raising an \code{AnchorError} on mismatch); the SM
uncertainty $\sigma_{\rm SM}=\sqrt{0.055^2+0.029^2}=0.062$~ps$^{-1}$ comes from the
HPQCD asymmetric uncertainties and is independently \emph{validated} against the
component quadrature in \code{\_validated\_sm\_theory\_sigma\_ps\_inv}. The band
also records a \code{central\_budget} ($|{\rm gap}|/2$) and a \code{tight\_budget}
($\max({\rm gap}-\sigma,\ \sigma_{\rm exp})/2$) for sensitivity studies; the HARD
veto uses the loose one.

The verdict: \code{evaluate\_bd\_mixing\_with\_running} returns
$\text{ratio} = |\Mtwelve^{\rm NP}|/\text{budget}_{\rm B001}$ and
\code{passes = (ratio <= 1.0)}. A point with $\text{ratio} > 1$ is
\textbf{vetoed}. The constraint class is
\begin{center}
\begin{tabular}{ll}
\toprule
\code{process\_id} & \code{"B001"} \\
\code{severity} & \code{Severity.HARD} \\
\code{observable} & \code{"Delta m\_d"} \\
required extra & \code{quark\_mass\_basis\_couplings} \\
\bottomrule
\end{tabular}
\end{center}

\section{The control flow, end to end}
\label{sec:codeflow}

\code{Constraint.evaluate(point)} does:
\begin{enumerate}[leftmargin=1.6em]
\item \code{couplings = point.get\_extra("quark\_mass\_basis\_couplings")}. If
\emph{absent}, it returns \code{passes=True} with a note that the constraint was
not evaluated (a graceful skip, not a silent pass of physics) and a diagnostic
\code{\{"missing\_extra": ...\}}.
\item \code{wilsons = bd\_mixing\_wilsons\_from\_couplings(couplings)} --- pulls
the \code{"b\_d"} Wilsons \eqref{eq:codeVLL}--\eqref{eq:codeLR}.
\item \code{result = bd\_mixing\_from\_wilsons\_with\_running(wilsons,
mu\_had=2.0, m12\_np\_budget=self.anchor.budget)} --- runs to $2$~GeV, contracts
the matrix elements, takes the magnitude, and divides by the B001 budget
\eqref{eq:b001budget}.
\item It also evaluates the \emph{core} budget result for the diagnostic
\code{core\_default\_ratio}/\code{core\_default\_passes}.
\item Returns a \code{ConstraintResult} with \code{predicted}
$=|\Mtwelve^{\rm NP}|$ (GeV), \code{ratio}, \code{budget}, the SM and experimental
central values, and a rich \code{diagnostics} dict
(\code{abs\_m12\_np\_gev}, \code{m\_kk\_gev}, \code{matching\_scale\_gev},
the complex \code{left\_db\_coupling}/\code{right\_db\_coupling}, the four
\code{wilson\_coefficients}, every budget variant, and full anchor provenance).
\end{enumerate}
The anchors are read from the sidecar \code{B001.yaml}: experiment from
\code{canonical\_experimental\_average}, SM from \code{standard\_model\_prediction}
(HPQCD~2019), and the GeV/bag constants from \code{auxiliary\_code\_inputs.\\deltaf2\_bd\_constants}, whose
\code{DELTA\_M\_BD\_EXP\_GeV = 3.334e-13} is the conversion anchor cross-checked
against the core.

% ====================================================================
\part{Subtleties and the gap to ``rigorous''}
\label{part:subtle}
% ====================================================================

\section{The operator basis is a subset (SLL/SRR absent)}
\label{sec:gap}
A pure vector KK-\emph{gluon} exchange at tree level generates VLL, VRR, and the
two LR operators \eqref{eq:VLL}--\eqref{eq:LR} --- but \emph{not} the scalar
$Q_2^{\rm SLL}/Q_3^{\rm SRR}$ operators of the full eight-operator basis. Blanke
et al.\ (0809.1073) carry the full set and note that for $B_{d,s}$ physics the
$Q_4^{\rm LR}$-type operator ``turns out to be as important as'' the others. Our
core is therefore a faithful tree-level KK-gluon matching, but it omits the
scalar operators that would arise from KK-Higgs or radion exchange and from
one-loop matching. For $\dmd$ this is a small effect (the dominant LR piece is
present), but it is an approximation, not the full EFT.

\section{Single-KK-mode truncation}
The matching keeps only the \emph{first} KK gluon (one $1/\MKK^2$ propagator). The
full tower sums to a slightly larger coefficient; the truncation is conservative
in the standard convention but is a known approximation. The same first-mode
choice fixes the $\MKK$-convention question below.

\section{The $\MKK$ convention}
\label{sec:Mconvention}
``$\MKK$'' is overloaded, exactly as in the oblique sector. The repo bookkeeping
default is \code{xi\_KK = 1.0}, i.e.\ $\MKK=\LIR$ (the geometric IR scale), per
\code{scales.py} (\code{DEFAULT\_QUARK\_XI\_KK = 1.0}). The \emph{physical} first
KK-gauge mass is $\MKK^{\rm phys}=x_1\LIR$ with the NN root
\code{GAUGE\_KK\_ROOT\_NN = 2.4487}, so a bound stated in $\LIR$ units is
$\approx2.45\times$ larger in physical units. When comparing our few-TeV $\dmd$
floor to a literature ``$\sim$~few TeV'' or to the CFW $21$/$33$~TeV context
values, the convention must match. The Wilson-coefficient reference scale
\code{DEFAULT\_QUARK\_TARGET\_SCALE\_GEV = 3000.0} (3~TeV) is a separate,
$\alpha_s$-anchoring scale and should not be conflated with $\MKK$.

\section{Leading-log QCD running and the bag-scale caveat}
The running is leading-log, and the bag parameters are FLAG-2024 values reported
at their own renormalisation scales; the kaon LR bags carry an explicit
$3~\text{GeV}\!\to\!2~\text{GeV}$ scale-matching \code{TODO} in the core. For
$B_d$ the $B_4/B_5$ inputs are used as stored; a fully NLO treatment would
RG-match the bags to $\mu_{\rm had}$ in the same scheme as the Wilson running.

\section{Coupling normalisation: $g_s$ vs $g_s^*$}
The default \code{compute\_quark\_kk\_gluon\_couplings} uses the perturbative
4D $g_s(\MKK)$ unless an explicit enhanced \code{g\_s\_star} is supplied
($\sim6$ at tree level, $\sim3$ at one loop, per the docstring citing CFW
0804.1954 and Gedalia--Grossman--Nir--Perez 2009). Whether a given scan uses the
plain or enhanced coupling shifts the floor by the corresponding factor; the
$\dmd$ \emph{magnitude} scales linearly in the product of the two off-diagonal
couplings, so a $g_s^*/g_s$ enhancement of $\sim2$ moves $|\Mtwelve^{\rm NP}|$ by
$\sim4$, i.e.\ the $\MKK$ floor by $\sim2$.

\section{Budget policy is a choice}
The HARD veto uses the \emph{loose} band \eqref{eq:b001budget}; the central and
tight variants are computed and stored but not used for the veto. This is a
deliberate, documented policy (an NP contribution as large as the
SM-vs-experiment gap \emph{plus} a combined-$\sigma$ cushion is tolerated). A
stricter analysis would use the central or tight band, tightening the $\dmd$
floor by up to a factor $\sim2.5$ in $|\Mtwelve^{\rm NP}|$ (so $\sim1.6$ in
$\MKK$). The ledger records this as a deliberate, uncertainty-aware choice rather
than an oversight.

\section{Checklist: what a stricter $\dmd$ treatment would add}
\begin{enumerate}[leftmargin=1.8em]
\item Add the scalar SLL/SRR operators (KK-Higgs/radion and one-loop matching).
\item Sum the full KK-gluon tower instead of the first mode.
\item Promote the leading-log running to NLO and RG-match the $B_d$ bag
parameters to $\mu_{\rm had}$ consistently.
\item Fix and document a single $\MKK$ convention (physical vs $\LIR$) across the
$\Delta F=2$ and electroweak sectors.
\item Decide the budget-band policy (loose / central / tight) on physics grounds
rather than convention.
\end{enumerate}

% ====================================================================
\part*{Common confusions}
\addcontentsline{toc}{section}{Common confusions}
% ====================================================================

\begin{callout}{Four common confusions, defused.}
\begin{itemize}[leftmargin=1.2em,nosep]
\item[\textbf{A.}] \emph{``$\dmd$ and $\epsK$ are the same constraint in two
mesons.''} No --- $\dmd$ reads the \textbf{magnitude} $|\Mtwelve|$ (CP-conserving),
$\epsK$ reads the \textbf{imaginary part} (CP-violating). The kaon also has a
$\sim$$200\times$ enhancement of its LR operator (bare chiral $r_\chi\!\approx\!26$
times LR running $\sim$$8$) versus a chiral prefactor $\sim$$1.6$ for $B_d$. Three
compounding factors make $\epsK$ the $\sim$$20$~TeV killer and $\dmd$ a few-TeV
constraint.
\item[\textbf{B.}] \emph{``Custodial symmetry relaxes the $\dmd$ bound.''} No ---
$P_{LR}$ protects the $Z\bar b_L b_L$ (and $Z$-mediated $\Delta F=2$) coupling,
but \code{B001} is a \textbf{KK-gluon} effect, blind to $P_{LR}$. What relaxes
$\dmd$ is \emph{flavour alignment} of the 5D couplings, a different lever.
\item[\textbf{C.}] \emph{``The budget is the SM--experiment gap.''} Partly --- the
\emph{core} budget is $\max(\dmd^{\rm exp}/2,\,|{\rm gap}|/2)=\dmd^{\rm exp}/2$,
but the \textbf{live HARD budget} \eqref{eq:b001budget} is the smaller
uncertainty-aware $(|{\rm gap}|+\sqrt{\sigma_{\rm exp}^2+\sigma_{\rm SM}^2})/2
\approx3.4\times10^{-14}$~GeV. Quote that one.
\item[\textbf{D.}] \emph{``$\MKK$ means one thing.''} It is overloaded by
$x_1\approx2.45$: the code default is $\MKK=\LIR$ (\code{xi\_KK=1.0}), while the
physical first KK mass is $2.45\,\LIR$. A few-TeV floor in $\LIR$ units is
$\sim2.45\times$ higher physically.
\end{itemize}
\end{callout}

% ====================================================================
\part*{Annotated reading guide}
\addcontentsline{toc}{section}{Annotated reading guide}
% ====================================================================

\begin{description}[leftmargin=0em,style=nextline]
\item[Sibling notes (same $\Delta F=2$ core).] The shared machinery --- operator
basis, KK-gluon matching, RG running, and hadronic matrix elements --- is stated
self-containedly above (\S\ref{sec:operators}, \S\ref{sec:rsmech},
\S\ref{sec:codewilsons}--\S\ref{sec:coderunning}) in the form the $B_d$
evaluation uses, so this note stands alone. The same core drives the kaon
($\epsK$, \code{epsilon\_k\_review.tex}), $B_s$ ($\Delta m_s$/$\phi_s$,
\code{delta\_m\_s\_review.tex}), and $D^0$ (\code{d0\_mixing\_review.tex})
reviews; the kaon note carries the fullest treatment of the LR chiral-plus-running
enhancement, and is the natural cross-read for the real-vs-imaginary contrast.

\item[Csaki, Falkowski, Weiler, arXiv:0804.1954.] The anarchic RS-flavour paper.
The tree-level KK-gluon $\Delta F=2$ Hamiltonian (their VLL/LR operators
\eqref{eq:VLL}--\eqref{eq:LR}), the chiral-plus-running enhancement of the kaon LR operator
($\sim$$200$ with the repo's bare $r_\chi\approx26$; CFW quote a combined
chiral-plus-running factor $\approx18$ in their own normalisation), the $g_s^*\sim6$ enhanced coupling, and the anarchic
$\MKK\gtrsim\mathcal{O}(20)$~TeV bound from $\epsK$ versus the much milder
$B_d$ bound. The $21$~TeV / $33$~TeV context values in our \code{B001.yaml}
\code{paper\_era\_reference} are theirs. This is the reference for
\S\ref{sec:operators}--\S\ref{sec:realvsimag}.

\item[Blanke, Buras, Duling, Gori, Weiler, arXiv:0809.1073.] The full
six/eight-operator $\Delta F=2$ analysis in custodially-protected warped models.
Read it for: the complete operator basis (including the SLL/SRR our core omits,
\S\ref{sec:gap}); the statement that $\epsK$ is the dominant constraint via the
LR chirality-flip operator; that $P_{LR}$ suppresses the $Z$-mediated piece but
does \emph{not directly relax} the $B_d$ KK-gluon bound (\S\ref{sec:custodial});
and the existence of ``modest fine-tuning'' / alignment regions reaching
$\MKK\simeq3$~TeV (\S\ref{sec:alignment}).

\item[HPQCD Collaboration, arXiv:1907.01025.] The SM lattice prediction
$\dmd^{\rm SM}=0.555\,{}^{+0.028}_{-0.055}(0.029)$~ps$^{-1}$ used as the theory
anchor \eqref{eq:dmdsm} and as the $\sigma_{\rm SM}$ in the B001 budget
\eqref{eq:b001budget}.

\item[FLAG Review 2024, arXiv:2411.04268.] Source of the $B_d$ decay constant and
bag parameters ($f_{B_d}, B_1, B_4, B_5$) used in the matrix elements
(\S\ref{sec:codeME}); auxiliary lattice context in the YAML.

\item[Belle II Collaboration, arXiv:2302.12791.] The modern direct $\dmd$
measurement ($0.516\pm0.008\pm0.005$~ps$^{-1}$, $190~\text{fb}^{-1}$) stored under
\code{recent\_experimental\_inputs}; not the live anchor, but a consistency
cross-check on \eqref{eq:dmdexp}.
\end{description}

\bigskip
\noindent\rule{\linewidth}{0.4pt}\\
\footnotesize
\emph{Internal note. The formulas \eqref{eq:codeVLL}--\eqref{eq:codeM12} and the
budgets \eqref{eq:corebudget}--\eqref{eq:b001budget} are exactly what
\code{quarkConstraints/deltaf2.py}, \code{physics\_adapters/deltaf2.py}, and
\code{B001.py} compute; the constants ($f_{B_d}=0.19$, $m_{B_d}=5.27972$,
$B_{1,4,5}=0.87/1.02/0.96$, \code{DELTA\_M\_BD\_EXP=3.334e-13},
\code{DELTA\_M\_BD\_SM=3.6e-13}, $\mu_{\rm had}=2$~GeV) are as stored in those
files and \code{B001.yaml}. ``Rigorous vs proxy'' tags follow the orchestration
ledgers: \code{B001} is \code{Severity.HARD} and rigorous. Cross-check before
citing externally.}

\end{document}
